% GENERATED FILE. Edit canonical lessons, then run npm run generate:books.
% canonical-source-hash: fnv1a64:63e7f396c2fc005b
% canonical-lessons: MR-C269-jhakan, MR-C269-barni, MR-C269-amti, MR-C269-usal, MR-C269-lonce

\chapter{Lid, Jar, Spiced lentil curry, Sprouted-bean curry, Pickle}
\label{ch:mr-lid-jar-spiced-lentil-curry-sprouted-bean-curry-pickle}
\hlchaptermodality{269}

\begin{chapteropening}
\textbf{By the end of this chapter:} \emph{I can say a lid, a jar, a spiced lentil curry, a sprouted-bean curry and a pickle.}

Complete the last lesson of chapter 269: I can say a lid, a jar, a spiced lentil curry, a sprouted-bean curry and a pickle.
\end{chapteropening}

\section[\texorpdfstring{\mr{झाकण} (jhākaṇ)}{झाकण (jhākaṇ)}]{\mr{झाकण} (jhākaṇ) — a lid}

\label{lesson:MR-C269-jhakan}



\begin{quote}
Before the new one: say the Marathi for a sitar, then the Marathi for a ball.
\end{quote}

\subsection*{You'll want to know: \mr{झाकण}}
\textbf{\mr{झाकण}} — \emph{jhākaṇ} — \textquotedblleft{}a lid\textquotedblright{}.

\begin{grammarlens}[title={What you've built}]
One more word for this chapter.
\end{grammarlens}

\subsection*{Your turn}
\begin{itemize}
\raggedright
  \item \emph{Say it:} \emph{jhākaṇ}
  \item \emph{Say it:} \emph{jhākaṇ}, once more
  \item \emph{Say it:} say \emph{jhākaṇ}
  \item \emph{Recall:} say the Marathi for a sitar, then the Marathi for a ball, then say \emph{jhākaṇ} again
\end{itemize}

\begin{tcolorbox}[breakable,colback=teal!4,colframe=teal!35!black,title={Before you move on}]
What does \mr{झाकण} mean? (\textquotedblleft{}A lid\textquotedblright{}.) Say it once more.
\end{tcolorbox}
\section[\texorpdfstring{\mr{बरणी} (barṇī)}{बरणी (barṇī)}]{\mr{बरणी} (barṇī) — a jar}

\label{lesson:MR-C269-barni}



\begin{quote}
Before the new one: say the Marathi for a ball, then the Marathi for a lid.
\end{quote}

\subsection*{You'll want to know: \mr{बरणी}}
\textbf{\mr{बरणी}} — \emph{barṇī} — \textquotedblleft{}a jar\textquotedblright{}.

\begin{grammarlens}[title={What you've built}]
One more word for this chapter.
\end{grammarlens}

\subsection*{Your turn}
\begin{itemize}
\raggedright
  \item \emph{Say it:} \emph{barṇī}
  \item \emph{Say it:} \emph{barṇī}, once more
  \item \emph{Say it:} say \emph{barṇī}
  \item \emph{Recall:} say the Marathi for a ball, then the Marathi for a lid, then say \emph{barṇī} again
\end{itemize}

\begin{tcolorbox}[breakable,colback=teal!4,colframe=teal!35!black,title={Before you move on}]
What does \mr{बरणी} mean? (\textquotedblleft{}A jar\textquotedblright{}.) Say it once more.
\end{tcolorbox}
\section[\texorpdfstring{\mr{आमटी} (āmṭī)}{आमटी (āmṭī)}]{\mr{आमटी} (āmṭī) — a spiced lentil curry}

\label{lesson:MR-C269-amti}



\begin{quote}
Before the new one: say the Marathi for a lid, then the Marathi for a jar.
\end{quote}

\subsection*{You'll want to know: \mr{आमटी}}
\textbf{\mr{आमटी}} — \emph{āmṭī} — \textquotedblleft{}a spiced lentil curry\textquotedblright{}.

\begin{grammarlens}[title={What you've built}]
One more word for this chapter.
\end{grammarlens}

\subsection*{Your turn}
\begin{itemize}
\raggedright
  \item \emph{Say it:} \emph{āmṭī}
  \item \emph{Say it:} \emph{āmṭī}, once more
  \item \emph{Say it:} say \emph{āmṭī}
  \item \emph{Recall:} say the Marathi for a lid, then the Marathi for a jar, then say \emph{āmṭī} again
\end{itemize}

\begin{tcolorbox}[breakable,colback=teal!4,colframe=teal!35!black,title={Before you move on}]
What does \mr{आमटी} mean? (\textquotedblleft{}A spiced lentil curry\textquotedblright{}.) Say it once more.
\end{tcolorbox}
\section[\texorpdfstring{\mr{उसळ} (usaḷ)}{उसळ (usaḷ)}]{\mr{उसळ} (usaḷ) — a sprouted-bean curry}

\label{lesson:MR-C269-usal}



\begin{quote}
Before the new one: say the Marathi for a jar, then the Marathi for a spiced lentil curry.
\end{quote}

\subsection*{You'll want to know: \mr{उसळ}}
\textbf{\mr{उसळ}} — \emph{usaḷ} — \textquotedblleft{}a sprouted-bean curry\textquotedblright{}.

\begin{grammarlens}[title={What you've built}]
One more word for this chapter.
\end{grammarlens}

\subsection*{Your turn}
\begin{itemize}
\raggedright
  \item \emph{Say it:} \emph{usaḷ}
  \item \emph{Say it:} \emph{usaḷ}, once more
  \item \emph{Say it:} say \emph{usaḷ}
  \item \emph{Recall:} say the Marathi for a jar, then the Marathi for a spiced lentil curry, then say \emph{usaḷ} again
\end{itemize}

\begin{tcolorbox}[breakable,colback=teal!4,colframe=teal!35!black,title={Before you move on}]
What does \mr{उसळ} mean? (\textquotedblleft{}A sprouted-bean curry\textquotedblright{}.) Say it once more.
\end{tcolorbox}
\section[\texorpdfstring{\mr{लोणचे} (loṇce)}{लोणचे (loṇce)}]{\mr{लोणचे} (loṇce) — a pickle}

\label{lesson:MR-C269-lonce}



\begin{quote}
Before the new one: say the Marathi for a spiced lentil curry, then the Marathi for a sprouted-bean curry.
\end{quote}

\subsection*{You'll want to know: \mr{लोणचे}}
\textbf{\mr{लोणचे}} — \emph{loṇce} — \textquotedblleft{}a pickle\textquotedblright{}.

\begin{grammarlens}[title={What you've built}]
One more word for this chapter.
\end{grammarlens}

\subsection*{Your turn}
\begin{itemize}
\raggedright
  \item \emph{Say it:} \emph{loṇce}
  \item \emph{Say it:} \emph{loṇce}, once more
  \item \emph{Say it:} say \emph{loṇce}
  \item \emph{Recall:} say the Marathi for a spiced lentil curry, then the Marathi for a sprouted-bean curry, then say \emph{loṇce} again
\end{itemize}

\begin{tcolorbox}[breakable,colback=teal!4,colframe=teal!35!black,title={Before you move on}]
What does \mr{लोणचे} mean? (\textquotedblleft{}A pickle\textquotedblright{}.) Say it once more.
\end{tcolorbox}
